%% P32 --- The transformer, derived
%%
%% PART IX'S CONTRACT IS THAT IT INTRODUCES NO NEW MATHEMATICS, so the first
%% question was not "what is left of the subject" but "which lines of the
%% brief have already been delivered".  `grep -rn 'prog:P32' programs/en/*.tex`
%% prints the deferrals with their contexts and it was the first thing run.
%%
%% EXACTLY THREE THINGS ARE THIS PROGRAM'S OWN:
%%   P03 and P06 both hand over the block's PARAMETER COUNT by name -- "the
%%       count itself is Program P32's, where every piece exists" -- and F02
%%       supplies the feed-forward half already, as 4d^2 + 4d^2.
%%   F12 hands over the RESIDUAL STREAM: "the architecture's own answer -- a
%%       path from the loss to the early layers that is not a long product at
%%       all -- is the residual connection, which Program P32 derives."  It
%%       names the mechanism and does not derive it; greps confirm nothing does.
%%   P25 hands over the SCORE MEASURED THROUGH THE ASSEMBLY: "E9 is measured
%%       here on random vectors rather than on a trained model.  Program P32
%%       assembles the architecture and measures it there."
%%
%% EVERYTHING ELSE IS ASSEMBLY AND IS SPENT: P05 the inner product, P25 the
%% sqrt(d_k) derived and elicited across its frames 30-41, P02 the stable
%% softmax, P19 the convex combination, P07 the reshape and the four axes,
%% P18 the layer-norm gradient, F08 the rotation, P03 the cache and the two
%% exponents, P04 the space.  None of it is re-taught.
%%
%% NOT SPENT HERE: whether the independence hypotheses survive TRAINING.  That
%% needs a trained model, which this book does not have -- so it is stated as
%% outstanding, on the precedent of P08, P11, P19, P20 and P25.
%%
%% Twin: programs/pl/P32-transformer-derived.tex.  Frame for frame, macro for
%% macro: parity.py compares the two in order, so a sentence rebuilt in Polish
%% has to carry the same maths spans and the same numeric literals in the same
%% sequence.  A digit stays a digit and a word stays a word.

\program{The transformer, derived}
\label{prog:P32}
\index{transformer}

\begin{outcomes}
\outcome{1--6}{name every operation in one block and say which program in
  this book already gave you it}
\outcome{7--12}{say why the score is divided by a square root, and say
  whether that argument survives the scores being produced rather than drawn}
\outcome{13--17}{say what the softmax buys the layer beyond making the
  numbers positive}
\outcome{18--20}{say why splitting a layer into heads costs nothing}
\outcome{21--25}{say why a residual connection changes what the gradient is
  a sum of, rather than merely making its factors larger}
\outcome{26--28}{say where a position enters a computation that is otherwise
  blind to order}
\outcome{29--35}{count a block's parameters, say which half of it holds
  more, and derive its cache from the same shapes}
\outcome{36--39}{say which of this book's derivations rest on a hypothesis
  that has not been checked on a trained model}
\end{outcomes}

\begin{quiz}
\item A block projects the stream four times and then runs a two-layer
  feed-forward network. Which of those two halves holds more parameters?
  \teachesat{29--30}
  \answerto{The feed-forward network, by a factor of $2$: attention is
  $4d^{2}$ and the feed-forward block is $8d^{2}$.}
\item Attention divides its scores by $\sqrt{d_k}$. In a block the queries
  and keys are not drawn, they are computed from the stream. Does the
  argument for that divisor still apply? \teachesat{8--9}
  \answerto{Yes, and it is exact: the assembled score is a bilinear form
  whose variance works out to $d_k$ in expectation, so nothing changes.}
\item A plain chain of $\val{p32.depth}$ layers has a gradient bounded by
  $\val{p32.plain.bound}$. What does adding a residual connection to each
  layer do to that? \teachesat{23--24}
  \answerto{It stops the gradient being a single product at all. It becomes
  a sum over paths, one of which is the identity and contributes $1$.}
\item A model with $\val{p32.layers}$ layers of width $\val{p32.d}$ is
  serving a conversation. Which grows faster with the sequence length, the
  arithmetic or the cache? \teachesatone{35}
  \answerto{The arithmetic, which is quadratic; the cache is linear. They
  are two exponents in one layer.}
\item Which of the claims this book has derived about attention has been
  checked on a model that was actually trained? \teachesat{36--37}
  \answerto{None of them. Every one is derived at initialisation, and this
  program says so rather than implying otherwise.}
\end{quiz}

% ============================================================
\section{One position through one block}
\label{sec:P32-one-block}
\index{transformer!one block}

\begin{fr}
This is the last part of the book, and it introduces no new mathematics.

That is not modesty about the material. It is the design: thirty-one
programs have each built one piece, and what is left is to put them in
order and see what the arrangement costs.

So here is a whole block, as six operations applied to one position's vector
in turn: project it three ways; take one dot product per pair of positions;
divide those by something; turn them into weights that sum to one; average
the third projection with those weights; then add the result back on to what
you started with and rescale.

Before reading on: how many of those six are operations this book has
already defined?
\dotline
\nextframe
\end{fr}

\begin{fr}
\ans{All six}

Every one, with the program that gave it to you:

\begin{itemize}
\item \textbf{Project it three ways.} A matrix is a function
  (Program~\ref{prog:P06}), and three of them applied to one vector is three
  functions applied to one argument.
\item \textbf{One dot product per pair.} An inner product
  (Program~\ref{prog:P05}), and Program~\ref{prog:F09} defined it before
  that.
\item \textbf{Divide by something.} Program~\ref{prog:P25} derives what,
  and why, and section~\ref{sec:P32-score} is about whether that derivation
  survives being assembled.
\item \textbf{Turn them into weights summing to one.} The $\softmax$, whose
  stable form is Program~\ref{prog:P02}'s.
\item \textbf{Average with those weights.} A convex combination
  (Program~\ref{prog:P19}).
\item \textbf{Add it back on and rescale.} The two halves of
  section~\ref{sec:P32-residual} and section~\ref{sec:P32-norm}.
\end{itemize}

There is nothing else in a block. What makes an architecture is the order
and the shapes, and the shapes are worth getting exactly right, because
Program~\ref{prog:P07} spent a program on what happens when they are not.

The stream carries a vector of width $\val{p32.d}$ at each position. The
queries are formed by a matrix $W_Q$. What shape must $W_Q$ be?
\dotline
\nextframe
\end{fr}

\begin{fr}
\ans{$\val{p32.d}$ by $d_k$}

It reads the stream, so it takes $\val{p32.d}$ numbers in; it produces a
query, so it puts $d_k$ out. Program~\ref{prog:P06} calls that a function
from one space to another and Program~\ref{prog:P07} calls it a type
signature; they are the same statement.

$W_K$ has the same shape, because a key has to meet a query in a dot
product and Program~\ref{prog:F09} requires two vectors of the same length
for that. $W_V$ takes the stream to whatever width the values are.

Three of the four projections are now accounted for: queries, keys and
values. The fourth reads the attention's output and writes into the stream.

Before reading on: why does that fourth one have to exist at all? The values
are already vectors \dash{} why not add the weighted average of them
straight back on?
\dotline
\nextframe
\end{fr}

\begin{fr}
\ans{Because the widths do not match, and because the heads have to be
combined}

The values live at width $d_k$ and the stream is $\val{p32.d}$ wide, so
something has to carry one to the other; that is Program~\ref{prog:P07}'s
type signature refusing an addition rather than a stylistic choice.

And once there are several heads, their outputs sit side by side and have to
be mixed. The fourth projection does both jobs at once, which is why it is
one matrix rather than two.

\textbf{Four projections, and that is the whole of attention's parameters.}
Hold on to the four; section~\ref{sec:P32-count} counts them.

Now the part that decides everything downstream. The scores are formed
between \emph{every} query position and \emph{every} key position. Over a
sequence of $n$ positions, how many scores is that?
\dotline
\nextframe
\end{fr}

\begin{fr}
\ans{$n^{2}$, and Program~\ref{prog:F10} counted them}

One per ordered pair. Program~\ref{prog:P03} priced that and made it a
warning: the parameters are $O(d^{2})$ and the work is $O(d\,n^{2})$, so a
count of parameters is not a count of work.

That $n^{2}$ is the single most consequential number in the architecture.
Every method with \enquote{long context} in its name is trying not to pay
it, and section~\ref{sec:P32-count} puts it beside the one quantity in a
block that is \emph{not} quadratic.

\begin{aibox}
Read a model card and you are given two numbers: a parameter count and a
context length. The first prices the matrices in this section and the second
prices the scores. They are different exponents in the same layer, so a card
that doubles its context has not doubled anything you can read off the
parameter count \dash{} it has quadrupled the score matrix and doubled the
cache. Neither figure appears on the card.
\end{aibox}
\end{fr}

\mermaidfig{p32-one-block}{One position through one block}{six operations}

\begin{fr}
One more shape before the arithmetic starts, because it is the one that
makes the rest of the program legible.

A block does not replace the stream. It \emph{adds} to it: the vector that
goes in comes out again with the block's output added on. So the stream runs
from the embedding at the bottom to the loss at the top, and every block
reads it, computes something, and adds the something back.

That is why it is called a stream rather than a pipeline, and
section~\ref{sec:P32-residual} shows that the choice of \emph{add} rather
than \emph{replace} is the reason a deep stack trains at all.
\end{fr}

% ============================================================
\section{The score, and whether the derivation survives assembly}
\label{sec:P32-score}
\index{attention!score}
\index{attention!the square-root divisor}

\begin{fr}
A score is one dot product: a query at one position against a key at
another. Program~\ref{prog:F09} defined it and Program~\ref{prog:P05} said
what it measures.

And it is divided by $\sqrt{d_k}$ before anything else happens to it.
Program~\ref{prog:P25} derived that divisor rather than asserting it, in a
section built around eliciting it, so you have produced this argument once
already.

Before reading on: say why the divisor is a square root, and why it is a
square root of the head width in particular.
\dotline
\nextframe
\end{fr}

\begin{fr}
\begin{ansblock}
A score is a sum of $d_k$ products. Variances of independent quantities add,
so a sum of $d_k$ of them has variance $d_k$ and spread $\sqrt{d_k}$.
Dividing by that root returns the score to a scale that does not depend on
$d_k$.
\end{ansblock}

Program~\ref{prog:P25} measured what happens without it: at
$d_k = \val{p25.e9.d.hi}$ the spread reaches $\val{p25.e9.raw.sd.512}$, one
key of $\val{p25.e9.keys}$ takes $\val{p25.e9.raw.top.512}$ per cent of the
weight, and the $\softmax$'s own derivative at that point is
$\val{p25.e9.resp.raw}$ against $\val{p25.e9.resp.scaled}$ with the division
in place.

So the divisor is a variance correction and not a constant somebody tuned.
None of that is this program's to prove again.

\textbf{But Program~\ref{prog:P25} measured it on vectors it drew directly},
and said so, and handed one thing here: in a block the queries and keys are
not drawn. They are \emph{computed}, both from the same stream, by two
matrices that a training run chose.

That is a different computation, and it is worth being precise about which
part of it is suspicious. Program~\ref{prog:P25}'s argument needs the entries
going into the dot product to be independent of one another. Here they are two
linear images of the \emph{same} stream, so they are functions of a shared
source, and \enquote{drawn independently} is exactly the hypothesis the
assembly might have destroyed.

It is worth knowing what turns on the answer before you give it. If the
hypothesis does not survive, then the divisor is correct about
Program~\ref{prog:P25}'s model and wrong about every block anybody has
trained, and this section would have to say so rather than repeating it. The
question is not a formality, and it is the only one in this program whose
answer the book had not already established somewhere.

Before reading on: does the variance argument still apply to a score computed
this way?
\dotline
\nextframe
\end{fr}

\begin{fr}
\ans{It does, and it is exact \dash{} but not for the reason it looks like}

Write the score out with the projections in place. The query is
$W_Q\T x$ and the key is $W_K\T y$, so
\[ q \cdot k \;=\; (W_Q\T x) \cdot (W_K\T y) \;=\; x\T M y,
   \qquad M = W_Q W_K\T. \]

That is worth stopping on. \textbf{The two projection matrices never appear
separately in a score.} Only their product does, and their product is one
$\val{p32.d}$-by-$\val{p32.d}$ matrix. A score is a bilinear form: one
vector, one matrix, one other vector.

And for two independent standard-normal streams the variance of $x\T M y$ is
exactly $\lVert M\rVert_F^{2}$ \dash{} the sum of the squares of every entry
of $M$. No sampling is needed to find it.

That is the whole of the assembly question, so it is worth naming the
detour it replaces.

\begin{note}
Two designs were tried before this one and both were wrong, and the second
failure is the interesting one.

The first drew fresh $W_Q$ and $W_K$ inside every trial. \textbf{That is not
a transformer.} A block has fixed weights and varying inputs, so averaging
over weight draws answers an easier question than the one that was asked.

The second fixed the weights and sampled the inputs, and was simply too
slow. Then the identity above made the sampling unnecessary altogether
\dash{} which is the finding rather than a shortcut, because a quantity you
can write down exactly does not need an error bar.
\end{note}
\end{fr}

\begin{fr}
Now the expectation, and it is four lines.

The fan-in rule Program~\ref{prog:P25} derives puts the variance of each
entry of a projection matrix at $1/\val{p32.d}$, so that a projected vector
comes out at the scale it went in. Then each entry of $M$ is a sum of $d_k$
products of two such entries, and

\[ \Ex\lVert M\rVert_F^{2}
   \;=\; \val{p32.d}^{2} \times d_k \times
         \Bigl(\frac{1}{\val{p32.d}}\Bigr)^{2}
   \;=\; d_k. \]

\textbf{The assembly preserves the variance exactly.} The two projections
between the stream and the score do not change the answer at all, so
Program~\ref{prog:P25}'s $\sqrt{d_k}$ is the right divisor for a real block
and not only for a pair of drawn vectors.

Measured, at $d_{\text{model}} = \val{p32.asm.dmodel}$ over
$\val{p32.asm.draws}$ weight draws, the assembled spread against the
prediction:

\begin{center}
\begin{tabular}{rrr}
\toprule
$d_k$ & assembled spread & $\sqrt{d_k}$ \\
\midrule
$8$  & $\val{p32.asm.sd.8}$  & $\num{2.83}$ \\
$16$ & $\val{p32.asm.sd.16}$ & $\num{4.00}$ \\
$32$ & $\val{p32.asm.sd.32}$ & $\num{5.66}$ \\
$64$ & $\val{p32.asm.sd.64}$ & $\num{8.00}$ \\
\bottomrule
\end{tabular}
\end{center}

The first and last rows are Program~\ref{prog:P25}'s own head sizes, and its
committed spreads there are $\val{p25.e9.raw.sd.8}$ and
$\val{p25.e9.raw.sd.64}$. The script asserts that this program's assembled
figures reproduce them, so the two programs cannot come apart about a
quantity they jointly describe.
\end{fr}

\begin{fr}
And now the half that Program~\ref{prog:P25}'s method could not reach, which
is the reason it handed the question over.

$\Ex\lVert M\rVert_F^{2} = d_k$ is a statement about the \emph{average} over
weight draws. A trained model does not have an average. It has one draw
\dash{} one $M$ \dash{} and what its scores actually do is
$\lVert M\rVert_F^{2}$ for that one matrix.

Measured across draws, the relative spread of that quantity:

\begin{center}
\begin{tabular}{rr}
\toprule
$d_k$ & spread of $\lVert M\rVert_F^{2}$ \\
\midrule
$8$  & $\val{p32.asm.spread.8}$\,\% \\
$16$ & $\val{p32.asm.spread.16}$\,\% \\
$32$ & $\val{p32.asm.spread.32}$\,\% \\
$64$ & $\val{p32.asm.spread.64}$\,\% \\
\bottomrule
\end{tabular}
\end{center}

\textbf{These are indicative rather than exact}, and the reason is worth
stating: $\val{p32.asm.draws}$ draws estimate a spread to about a seventh of
itself, so the column is a size and not a measurement. What is not
indicative is the ordering, and the ordering has a reason.

It is worth being clear what a spread across draws means for something you
deploy. Two models trained from different seeds begin from different $M$, so
their scores sit on slightly different scales before either has seen any
data \dash{} which is one of the places seed-to-seed variation in a training
run comes from, rather than anything the optimiser did.

Before reading on: $M$ has $\val{p32.d}^{2}$ entries whatever $d_k$ is. Why
should a \emph{narrow} head give a more variable $\lVert M\rVert_F^{2}$ than
a wide one?
\dotline
\nextframe
\end{fr}

\begin{fr}
\ans{Because the rank of $M$ is at most $d_k$}

$M = W_Q W_K\T$ is a product through a $d_k$-wide waist, so
Program~\ref{prog:P08}'s bottleneck theorem applies to it directly: its rank
cannot exceed $d_k$ however large it is.

So the entries of $M$ are not $\val{p32.d}^{2}$ independent numbers. They are
$\val{p32.d}^{2}$ numbers built out of $d_k$ independent directions, and a
narrow head has fewer of them to average over. That is why the spread falls
as the head widens, and it is Program~\ref{prog:P08} doing a job nobody
wrote it for.

\begin{trapbox}
\textbf{\enquote{The scaling is exact, so a real model's scores have
variance $d_k$.}}

The first clause is right and the second does not follow. Exact in
expectation is a statement about the ensemble of models you might have
trained; a deployment has one, and its scores are off by a few per cent
before training starts.

That is Program~\ref{prog:P21}'s distinction \dash{} unbiased is not the
same as usable \dash{} arriving in an architecture instead of an estimator.
It is a small effect here and it is worth knowing it is not zero, because
the next section's question is what happens to it after training, and
nothing in this book answers that.
\end{trapbox}
\end{fr}

\mermaidfig{p32-score-assembled}{Where a score comes from}{one bilinear form}

% ============================================================
\section{From scores to weights}
\label{sec:P32-softmax}
\index{transformer!softmax over scores}

\begin{fr}
The scores are now a row of real numbers, one per key position, and they can
be any sign and any size. They have to become weights.

Program~\ref{prog:F05} settled what the $\softmax$ does to them and
Program~\ref{prog:P02} settled how to compute it without overflowing. Both
are spent; neither is repeated here.

What is worth asking is what the layer \emph{buys} by that choice, because
positivity alone would be much cheaper. Squaring every score gives positive
numbers; so does exponentiating and stopping there. Neither of those sums to
one, and the division at the end of a $\softmax$ is the only step that makes
it do so \dash{} so the question is not what the exponential buys but what
the \emph{normalisation} buys.

The weights come out positive and summing to one. Given that, what can you say
about where the layer's output must lie, relative to the value vectors it is
averaging?
\dotline
\nextframe
\end{fr}

\begin{fr}
\ans{Inside their convex hull}

That is Program~\ref{prog:P19}'s object exactly. Weights that are
non-negative and sum to one make a convex combination, so the output is a
weighted average of the values and can never leave the region they span.

It is a real guarantee and it is worth naming, because it says what a single
attention layer structurally \emph{cannot} do: it cannot produce a direction
that none of its values point in. Whatever it writes into the stream is
built from what was already there.

The feed-forward block that follows is what escapes that, because
Program~\ref{prog:F05} proved a composition of affine maps is affine and a
non-linearity is what breaks it. The two halves of a block do different
jobs, and this is the sentence that says which.

\begin{aibox}
The convex-hull property is why a purely attentional stack cannot invent a
feature: at every layer the output sits among its own inputs. When an
interpretability result reports a direction that appears at layer $k$ and
not at layer $k-1$, the feed-forward block is where it can have come from,
and Program~\ref{prog:P31} is the program about what such a result does and
does not license.
\end{aibox}
\end{fr}

\begin{fr}
There is a cheaper way to make numbers sum to one, and it is worth seeing
why it is not used: divide each score by the total of the scores.

Before reading on: the scores at this point can be any sign. What does that
do to the cheap method?
\dotline
\nextframe
\end{fr}

\begin{fr}
\ans{It breaks it, in two separate ways}

A negative score would give a negative weight, and the output would no longer
be an average of the values at all \dash{} the convex-hull guarantee is gone.
Worse, the scores can sum to zero, and then the method divides by zero for
reasons that have nothing to do with the data being unusual.

Exponentiating first fixes both: it sends every score to a positive number,
and a sum of positive numbers is never zero. \textbf{The exponential is
there to make the division safe and the guarantee true}, and the fact that it
also turns differences of scores into ratios of weights is what
Program~\ref{prog:F07} showed and what makes the divisor of
section~\ref{sec:P32-score} matter so much.
\end{fr}

\begin{fr}
One caution about the weights, and it is Program~\ref{prog:P29}'s.

A weight of $\num{0.9}$ on one key is often read as the layer \enquote{attending
to} that key, and the reading is a statement about a distribution, so it can
be quantified rather than eyeballed: Program~\ref{prog:P29} converts an
attention row's entropy into an effective number of keys, and reports
$\val{p29.att.raw.eff}$ of $\val{p29.att.keys}$ without the divisor against
$\val{p29.att.scaled.eff}$ with it.

So the divisor of section~\ref{sec:P32-score} is not a numerical nicety.
It is what stops the layer collapsing onto one key before it has learnt
anything, and its effect is readable in the unit a design review argues in.
\end{fr}

% ============================================================
\section{Heads}
\label{sec:P32-heads}
\index{transformer!multi-head}

\begin{fr}
A block does not run one attention, it runs several in parallel and
concatenates the results. Program~\ref{prog:P07} owns the mechanics
\dash{} the wide vector is split into $h$ slices, the head axis moves in
front of the position axis, one contraction does the rest \dash{} and calls
it a reshape rather than a computation.

Here is what that costs. With $\val{p32.heads}$ heads of width
$\val{p32.d.head}$, the head widths add to
$\val{p32.heads} \times \val{p32.d.head} = \val{p32.d}$, which is the
stream's own width.

Before reading on: compare the parameters in $\val{p32.heads}$ heads of
width $\val{p32.d.head}$ against one head of width $\val{p32.d}$.
\dotline
\nextframe
\end{fr}

\begin{fr}
\ans{They are the same, exactly}

$W_Q$ for one wide head is $\val{p32.d}$ by $\val{p32.d}$. For
$\val{p32.heads}$ narrow heads it is $\val{p32.heads}$ matrices of
$\val{p32.d}$ by $\val{p32.d.head}$ \dash{} and
$\val{p32.heads} \times \val{p32.d.head} = \val{p32.d}$, so the columns
add up to the same matrix.

\textbf{So multi-head attention is free.} It is one $\val{p32.d}$-by-$\val{p32.d}$
projection read as $\val{p32.heads}$ blocks rather than one, which is why
Program~\ref{prog:P07} could call it a reshape: nothing is computed
differently, the numbers are grouped differently.

What changes is not the cost but the constraint. One wide head forms one
$M = W_Q W_K\T$ of rank up to $\val{p32.d}$; $\val{p32.heads}$ narrow heads
form $\val{p32.heads}$ separate ones, each of rank at most
$\val{p32.d.head}$. The block trades one rich comparison for many restricted
ones at identical cost \dash{} and Program~\ref{prog:P08} is where the word
\enquote{restricted} was made precise.

Program~\ref{prog:P07} makes one point about the arrangement that is easy
to skip. The head axis is moved \emph{in front of} the position axis before
the contraction.

Before reading on: what would go wrong if it were not?
\dotline
\nextframe
\end{fr}

\begin{fr}
\ans{Positions from different heads would meet}

The contraction runs over the last axis and pairs everything in front of it
independently. With the head axis in front, every head compares its own
queries against its own keys and the heads never mix. With the axes the other
way round, the same contraction would pair a query from one head against a
key from another, which is not a comparison of anything.

No error would be raised. The shapes work out either way, and
Program~\ref{prog:P07} is a whole program about that being the dangerous
case \dash{} an arrangement that runs and computes something other than what
was meant.

\begin{trapbox}
\textbf{\enquote{More heads means more capacity.}}

More heads at fixed width would. More heads at fixed \emph{total} width
\dash{} which is what every architecture actually does \dash{} means the
same parameters carved into narrower, lower-rank pieces. Whether that trade
is worth it is an empirical question about a class of task, and this book
does not answer it; what the arithmetic settles is that it is a trade rather
than a gain.
\end{trapbox}
\end{fr}

% ============================================================
\section{The residual stream}
\label{sec:P32-residual}
\index{transformer!residual stream}
\index{gradient!and residual connections}

\begin{fr}
Program~\ref{prog:F12} proved something about depth and then declined to
finish it.

Its argument: the gradient through a stack of layers is a \emph{product},
one factor per layer, and if each factor is at most $\frac{1}{4}$ then the
whole thing is at most $\bigl(\frac{1}{4}\bigr)^{n}$.

Before reading on: at $\val{p32.depth}$ layers, what is that bound?
\dotline
\nextframe
\end{fr}

\begin{fr}
\ans{$\val{p32.plain.bound}$, and that is the best case}

Program~\ref{prog:F12} commits exactly that number and this program
recomputes it on the same chain, so if either moves the build stops.

It is worth being clear how bad it is, and also how it is not bad in quite
the way people say. At $\val{p32.plain.bound}$ the number is still perfectly
representable: Program~\ref{prog:P01} puts $\code{float32}$'s smallest normal
at $\val{p01.fp32.minnorm}$, so nothing underflows in the best case.

What underflows is the realistic one. Program~\ref{prog:F12} evaluates the
same product at a \emph{saturated} layer rather than at its best, and gets
$\val{f12.sat.bound}$ \dash{} below every floor that format has. Then the
early layers receive exactly zero, and receive it silently.

Program~\ref{prog:F12} then names the architecture's answer and hands the
derivation here: \emph{a path from the loss to the early layers that is not
a long product at all.}

So here is the change, and it is one symbol. Instead of a layer replacing
the stream, $y = f(x)$, it adds to it:
\[ y \;=\; x + f(x). \]

Two things about that are worth noticing before you differentiate it. It puts
a constraint on the block: $x$ and $f(x)$ have to have the same shape to be
added at all, which is why every width in
section~\ref{sec:P32-one-block} came back to $\val{p32.d}$. And it introduces
no operation this book has not got \dash{} it is an addition, and
Program~\ref{prog:F11} gives the derivative of a sum.

So nothing here needs a new rule, and a change that needs no new rule is
exactly the kind that gets read as cosmetic. It is worth deciding for yourself
what it does before being told what it is for.

Before reading on: differentiate that. What is the Jacobian of one such
layer?
\dotline
\nextframe
\end{fr}

\begin{fr}
\ans{$1 + f'(x)$}

Which looks like a small change and is not. Stack $n$ of them and the
gradient is
\[ \prod_{k=1}^{n} \bigl(1 + f'_k\bigr), \]
where before it was $\prod_k f'_k$. The obvious reading is that the factors
have been made bigger, and that reading is available and wrong.

Multiply the brackets out. Every term picks, from each bracket, either the
$1$ or the $f'_k$ \dash{} so the expansion has one term per \emph{subset} of
the layers, $2^{n}$ of them, and each term is the product of the $f'$ over
the layers in its subset.

\textbf{That is a sum over paths.} A subset says which layers the gradient
went \emph{through} and which it went \emph{around}.

Notice what the expansion has already done to the question. A product is a
single quantity, and one small factor anywhere in it settles the whole thing.
A sum is $2^{n}$ quantities, and it is settled by whichever of them is
largest, so a small term costs nothing. The two expressions are equal, and
they are not equally easy to reason about \dash{} which is why multiplying
out was worth the line.

Before reading on: one of those $2^{n}$ subsets is the empty one. What does
its term contribute?
\dotline
\nextframe
\end{fr}

\begin{fr}
\ans{Exactly $1$}

The empty subset takes the $1$ from every bracket. It is the path that goes
around every layer \dash{} straight from the loss to the input, touching
nothing.

And that is the whole thing:

\begin{center}
\textbf{a plain chain's gradient is one product, which $n$ small factors
kill.\\
A residual stack's is a sum, and one term of that sum is $1$.}
\end{center}

No factor can kill it, because it does not contain any factors. The other
$2^{n} - 1$ paths can each vanish exactly as Program~\ref{prog:F12} says
they will, and the gradient still arrives.

The script checks the expansion over fractions at
$\val{p32.path.n}$ layers \dash{} $\val{p32.path.count}$ paths, summing to
the product exactly \dash{} and then checks the consequence, which is the
claim rather than the arithmetic: turn every layer off, so that every
$f'_k$ is zero, and a plain chain's gradient is exactly $0$ while a residual
stack's is exactly $1$.
\end{fr}

\transcript{p32-two-stacks}

\begin{fr}
The second line of that listing is the honest case, where the layers do
something small rather than nothing.

Over $\val{p32.res.trials}$ draws with each $f'$ of size about
$\val{p32.res.sd}$ and $\val{p32.depth}$ layers, the plain chain's median
gradient is $\val{p32.res.plain.med}$ and the residual stack's is
$\val{p32.res.resid.med}$, with $\val{p32.res.within}$ per cent of draws
within a factor of two of $1$.

Those two numbers are an illustration of an exact result and not the result,
so what the script asserts is the ordering and the fact that the residual
stack stays of order one, never the figures themselves.

\begin{trapbox}
\textbf{\enquote{Residual connections fix vanishing gradients.}}

They change what the gradient \emph{is}. That is stronger than fixing it and
it is also narrower, and both halves matter.

Stronger: the identity path is not a large factor that might yet be
outweighed, it is a term no factor multiplies.

Narrower: nothing here says the other $2^{n} - 1$ paths behave. A residual
stack can still explode \dash{} Program~\ref{prog:F12}'s second failure mode
is untouched, which is why the next section exists \dash{} and adding
residual connections to a badly scaled network gives you a gradient of $1$
plus a great deal of noise. The identity path guarantees arrival, not
usefulness.
\end{trapbox}
\end{fr}

\mermaidfig{p32-paths-not-product}{Two stacks, two gradients}{sum, not product}

% ============================================================
\section{Normalisation, and where a position enters}
\label{sec:P32-norm}
\index{transformer!layer normalisation}
\index{transformer!positional information}

\begin{fr}
Two operations are left and both are spent, so this section says where they
sit rather than what they are.

\textbf{Normalisation} subtracts the mean of a vector and divides by its
spread. Program~\ref{prog:P18} derives its gradient and checks it against a
finite difference to better than $\val{p18.check.ceiling}$, so nothing about
it is owed here.

What the assembly adds is why it is in the block at all. Section
\ref{sec:P32-residual} left the second failure mode open: the identity path
guarantees the gradient arrives, and says nothing about the size of what
arrives with it. Every block adds to the stream, so the stream's scale grows
with depth unless something holds it.

Before reading on: what does normalising do to the argument of
section~\ref{sec:P32-score}?
\dotline
\nextframe
\end{fr}

\begin{fr}
\ans{It restores the condition the argument assumes}

The variance calculation assumed the stream's entries had a known, fixed
scale. Thirty-two blocks adding into it would not leave that true. A
normalisation puts it back before each block reads it, so the score's
variance argument holds at layer $\val{p32.layers}$ for the same reason it
holds at layer one.

So the three pieces are one mechanism rather than three: the residual path
carries the gradient, the normalisation keeps the scale it is carried at
fixed, and the divisor keeps the scores on that scale.

\textbf{Positional information} is the second, and it repairs something
the block as built so far simply does not have.

Before reading on: take the block of section~\ref{sec:P32-one-block} and
shuffle the input positions. What happens to its outputs?
\dotline
\nextframe
\end{fr}

\begin{fr}
\ans{They shuffle with it, and are otherwise identical}

Every score is a dot product between two vectors and the weighted sum runs
over positions without consulting them, so nothing in the block can tell
\enquote{the cat sat} from \enquote{sat cat the}. It is not that order is
handled badly; it is that order is not represented at all.

Program~\ref{prog:F08} supplies the repair and proves the property it turns
on: rotating a query and a key by angles $\alpha$ and $\beta$ leaves their
dot product depending only on $\beta - \alpha$. So rotating each position's
vector by an angle proportional to its index makes every score a function of
the \emph{distance} between two positions, and of nothing else about where
they sit.

The rotation is orthogonal, so it changes no length \dash{}
Program~\ref{prog:F08} measures the error at $\val{f08.rot.len.err}$ and
Program~\ref{prog:P09} says why an orthogonal change of basis costs the
model nothing. \textbf{Position is added without spending any of the
stream's capacity}, which is the whole reason the method is used.
\end{fr}

% ============================================================
\section{What it costs}
\label{sec:P32-count}
\index{transformer!parameter count}
\index{cost!of one block}

\begin{fr}
Program~\ref{prog:F01} opened this book by having you work out what a
seven-billion-parameter model weighs. It never said where the seven billion
comes from, and on page three there was no vocabulary for it.

There is now. Program~\ref{prog:P03} commits the shape of the same model
\dash{} $\val{p32.layers}$ layers of width $\val{p32.d}$ \dash{} and
section~\ref{sec:P32-one-block} counted the matrices.

Attention has four projections, each $\val{p32.d}$ by $\val{p32.d}$ once the
heads are put back together, so it holds $4d^{2}$.

The feed-forward block is two layers, $d$ in and $4d$ out, then $4d$ in and
$d$ out. Program~\ref{prog:F02} collected those like terms already:
$4d^{2} + 4d^{2} = 8d^{2}$.

Both counts are per block, so neither of them moves with depth: multiplying
by the number of layers multiplies both by the same thing and leaves their
ratio exactly where it is. And neither is new \dash{} the first is four
projections of one width and the second is
Program~\ref{prog:F02}'s collected like terms, so the comparison is between
two numbers this book has already handed you.

Both also count matrices and nothing else. A block carries biases and two
sets of normalisation parameters as well, and those are $\val{p32.d}$ numbers
apiece against matrices of $\val{p32.d}^{2}$ \dash{} smaller by a factor of
the width, which is why nobody counts them and why it is worth saying once
that they were left out rather than overlooked.

Before reading on: which of the two halves of a block holds more parameters,
and by how much?
\dotline
\nextframe
\end{fr}

\begin{fr}
\ans{The feed-forward block, by a factor of $2$}

$8d^{2}$ against $4d^{2}$. \textbf{The half nobody discusses holds twice
what the half everybody discusses does.}

It is worth sitting with, because almost every sentence written about this
architecture is about attention, and two thirds of a block's parameters are
in the part that is not attention at all. In numbers, at
$d = \val{p32.d}$:

\begin{center}
\begin{tabular}{lr}
\toprule
attention, $4d^{2}$ & $\val{p32.attn.block}$ \\
feed-forward, $8d^{2}$ & $\val{p32.ffn.block}$ \\
\midrule
one block, $12d^{2}$ & $\val{p32.per.block}$ \\
$\times\ \val{p32.layers}$ layers & $\val{p32.blocks.total}$ \\
\bottomrule
\end{tabular}
\end{center}

And now the arithmetic Program~\ref{prog:F01} could not do.

Note what the two numbers are. The one below is a model card's figure, which
is what anybody would look up; the one above is what the shapes force, and
nobody chose it. So whatever gap there is between them has to be accounted for
by something this count has left out, and it is worth deciding whether you
expect that to be a rounding or a component.

Before reading on: $\val{p32.blocks.total}$ against
Program~\ref{prog:F01}'s $\val{p32.f01.params}$ \dash{} what fraction of
that model is the blocks?
\dotline
\nextframe
\end{fr}

\begin{fr}
\ans{$\val{p32.blocks.pct}$ per cent of it}

\textbf{So the blocks are the model.} The remainder is the embedding, which
at a vocabulary of $\val{p32.vocab}$ is $\val{p32.embed}$ parameters, or
$\val{p32.embed.pct}$ per cent \dash{} and the rest is the output head, the
biases and the normalisation scales, which are counted in thousands where
these are counted in billions.

The book's first program had you weigh a number it could not derive. Thirty
programs later the number comes out of two committed shapes and one
multiplication.
\end{fr}

\begin{fr}
Two honest caveats about that count, because it is the kind of number that
travels.

\begin{warning}
\textbf{The $12d^{2}$ is the original block's, not every block's.} The
feed-forward half is $8d^{2}$ only when its inner width is $4d$ and it uses
two matrices. Several current architectures use three matrices and an inner
width that is not $4d$, which changes the $8$ and therefore the $12$. The
method survives \dash{} count the matrices, add the squares \dash{} and the
constant does not.

And the vocabulary is a \emph{design choice} rather than a measurement.
Program~\ref{prog:F01} never fixed one, so $\val{p32.vocab}$ is named here
as the assumption it is; a model with four times the vocabulary moves the
embedding row and not the block rows.
\end{warning}
\end{fr}

\begin{fr}
Now the cache, and it comes out of the same shapes.

Serving a conversation means not recomputing the keys and values of every
earlier position at every step, so they are kept. What has to be kept is one
key and one value per position, per layer.

Notice which of the four projections that is, and which it is not. The
queries are not kept: at each step there is exactly one new query, it is used
once against every key, and it is never wanted again. The keys and the values
are the ones every future step will need. \textbf{Two of the four are stored
and two are not}, and which two follows from what the computation does with
them rather than from a choice anybody made.

Before reading on: at $\val{p32.layers}$ layers of width $\val{p32.d}$, two
bytes a number, how much is that per token?
\dotline
\nextframe
\end{fr}

\begin{fr}
\ans{$2 \times \val{p32.layers} \times \val{p32.d} \times \val{p32.bytes}$
bytes, which is $\val{p32.kv.per.token.mib}$ MiB}

Two tensors, one $d$-vector each, per layer, per position. Over
$\val{p32.seq}$ tokens that is $\val{p32.kv.gib}$ GiB \dash{} which is
Program~\ref{prog:P03}'s own committed figure, and the script asserts that
this program's derivation from the block's shapes reproduces it.

That is the assembly earning its place. Program~\ref{prog:P03} measured the
cache; it could not say why it had that size, because saying so needs the
four projections and the fact that exactly two of them are stored. Now it
does.

\begin{note}
Program~\ref{prog:P03} printed this per-token figure as \enquote{MB} while
computing it in mebibytes, and it survived a whole draft because
$\val{p32.kv.per.token.mib}$ MiB and $\num{0.524}$ MB both round to
$\val{p32.kv.per.token.mib}$ at one decimal. The unit is corrected there
now. It is the exact confusion Program~\ref{prog:P03}'s own summary warns
about \dash{} a capacity plan that is seven per cent out is nearly always
this one \dash{} and the reason it hid is Program~\ref{prog:P17}'s: two
readings of a formula that agree numerically stay invisible until they do
not.
\end{note}
\end{fr}

\begin{fr}
And the last arithmetic in the program is the one that decides how the
architecture is deployed.

Over $\val{p32.seq}$ positions with $\val{p32.heads}$ heads, the scores are
$\val{p32.scores}$ numbers. The cache is $\val{p32.cache.vectors}$ vectors.

Double the sequence and the first goes up \emph{fourfold} while the second
goes up \emph{twofold} \dash{} two exponents in one layer, which is
Program~\ref{prog:P03}'s finding, now read off the shapes rather than
asserted about attention in general.

The quadratic one is the thing to hold on to. \textbf{Every method with
\enquote{long context} in its name is an attempt not to pay $n^{2}$}:
computing only some of the pairs, computing them at lower precision,
computing them in blocks that fit in fast memory, or replacing the
comparison with something that is not a comparison of every pair. They
differ in what they give up, and what they give up is which pairs are
allowed to see each other.

\begin{aibox}
The two exponents explain a pattern anybody who has served this
architecture has met. A model that answers a short prompt comfortably and
falls over on a long one has usually not run out of \emph{parameters}: it has
run out of one of these two. Which one tells you what to do, and they point
opposite ways \dash{} the cache is fixed by serving fewer conversations at
once, the score matrix by serving shorter ones. A capacity plan that does not
separate them is guessing.
\end{aibox}
\end{fr}

\mermaidfig{p32-where-parameters-are}{Where a block's parameters sit}{four and eight}

% ============================================================
\section{What this book has not checked}
\label{sec:P32-unchecked}
\index{measurement!what is outstanding}

\begin{fr}
The architecture is assembled and every number in it came from a program
that derived it. This last section is about the ones that were derived
somewhere the model does not live.

Section~\ref{sec:P32-score} proved that the assembly preserves the score's
variance. Read the proof again and ask what it assumed about the stream.
\dotline
\nextframe
\end{fr}

\begin{fr}
\ans{That its entries were independent, and all of one scale}

Which is true at initialisation, because the initialiser makes it true.
Program~\ref{prog:P25} says so in as many words and adds the part that
matters: it becomes less true as training proceeds, because queries and keys
learn structure, and structure is dependence.

So every one of this book's derivations about attention \dash{} the
$\sqrt{d_k}$, the variance through the projections, the fan-in rule, the
scale a normalisation restores \dash{} is a statement about a model
\textbf{before it has been trained}.

\begin{warning}
\textbf{This book has not measured any of them on a trained model, and it
does not have one.}

The honest position is that the derivations are exact and their hypothesis
is checkable only where this book cannot go. What would settle it is not
difficult and it is not something a sandbox can do: take a trained model,
compute $\lVert W_Q W_K\T\rVert_F^{2}$ per head, and compare it against
$d_k$. If the two have parted company, the divisor is no longer doing what
it was derived to do, and every downstream claim about the entropy of an
attention row is about a distribution nobody has checked the scale of.

The prediction the mathematics makes is that they part company \emph{and it
does not matter}, because $\sqrt{d_k}$ stays the right order of magnitude
whatever the dependence does. That prediction is a judgement. It is
labelled as one here rather than presented as a result.
\end{warning}
\end{fr}

\begin{fr}
It is not the only one, and the list is worth having in one place because
each entry names a program that stated its own limit rather than papering
over it.

\begin{center}
\begin{tabularx}{\linewidth}{lX}
\toprule
Program & What is not measured \\
\midrule
\ref{prog:P08}, \ref{prog:P11} & whether a real embedding matrix's spectrum
  decays as the low-rank argument assumes \\
\ref{prog:P19} & whether which basin a walk lands in matters at the scale
  people train at \\
\ref{prog:P20} & whether an adaptive method reaches a better answer, rather
  than reaching one faster \\
\ref{prog:P25}, \ref{prog:P32} & whether the independence the score's
  derivation assumes survives training \\
\ref{prog:P31} & what a probe's score licenses about a layer, once the
  gap between the bound and the truth is measured \\
\bottomrule
\end{tabularx}
\end{center}

\textbf{Five entries and one cause.} Every one needs a trained model's real
matrices, and this book was written where there is not one. That is a
limitation of the book rather than of the mathematics, and the reason it is
printed here rather than left implicit is the discipline both companion
volumes are built on: a claim needs a method and a number, or it is labelled
judgement.
\end{fr}

\begin{fr}
One closing observation, and it is about what the assembly showed rather
than about the architecture.

Nothing in this program was new. The score is Program~\ref{prog:P05}'s inner
product; its divisor is Program~\ref{prog:P25}'s variance argument; the
weights are Program~\ref{prog:P19}'s convex combination; the heads are
Program~\ref{prog:P07}'s reshape; the residual path is
Program~\ref{prog:F12}'s product read as a sum; the count is
Program~\ref{prog:F02}'s collection of like terms done twelve times; the
cache is Program~\ref{prog:P03}'s measurement derived from shapes.

And yet three things came out of putting them together that none of those
programs could have said alone: that the projections leave the variance
argument \emph{exactly} intact, that a residual connection changes what the
gradient \emph{is} rather than how big it is, and that the feed-forward
block holds twice what attention does.

\textbf{That is what an architecture is.} Not a new idea, but an order in
which known ideas are applied, chosen so that each one's assumptions are
still true when the next one needs them. Reading it that way is the thing
this book was for.
\end{fr}

\begin{summarybox}
\sumitem{1--2}{A block is six operations and this book had already defined
  all six. What makes an architecture is the order and the shapes.}
\sumitem{3--5}{Attention's parameters are four projections; a score is one
  inner product; and $n$ positions give $n^{2}$ scores, which is the most
  consequential number in the design.}
\sumitem{7--8}{The score's divisor is Program~\ref{prog:P25}'s variance
  correction. In a block the queries and keys are \emph{produced}, so the
  question is whether that argument survives, and Program~\ref{prog:P25}
  handed it here.}
\sumitem{9--10}{It does, exactly. A score is a bilinear form
  $x\T M y$ with $M = W_Q W_K\T$, its variance is $\lVert M\rVert_F^{2}$,
  and that has expectation $d_k$ \dash{} so the projections change nothing
  and no sampling was needed to find out.}
\sumitem{11--12}{But a trained model has one weight draw and not an average,
  so its own scores are a few per cent off before training starts, and the
  spread falls as the head widens because $M$'s rank is at most $d_k$
  (Program~\ref{prog:P08}).}
\sumitem{13--16}{The $\softmax$ makes the output a convex combination
  (Program~\ref{prog:P19}), so a purely attentional layer cannot produce a
  direction its values do not already span.}
\sumitem{18--20}{Multi-head attention is free: $\val{p32.heads}$ heads of
  width $\val{p32.d.head}$ hold exactly the parameters of one head of width
  $\val{p32.d}$. What changes is rank, not cost.}
\sumitem{21--22}{A plain chain of $\val{p32.depth}$ layers has a gradient of
  at most $\val{p32.plain.bound}$, which is below what a $\code{float32}$
  can hold.}
\sumitem{23--25}{$y = x + f(x)$ makes the Jacobian $\prod(1 + f'_k)$, whose
  expansion is a sum over $2^{n}$ paths \dash{} one per subset of layers.
  \textbf{The empty subset contributes exactly $1$.} A plain chain's
  gradient is one product; a residual stack's is a sum with a $1$ in it.}
\sumitem{26--27}{Normalisation restores the scale the score's argument
  assumes, which is why the three pieces are one mechanism.}
\sumitem{28}{A rotation makes a score depend on the distance between two
  positions and nothing else, and costs the model nothing because it is
  orthogonal (Programs~\ref{prog:F08} and \ref{prog:P09}).}
\sumitem{29--32}{Attention is $4d^{2}$ and the feed-forward block is
  $8d^{2}$: \textbf{the half nobody discusses holds twice the half everybody
  does}. A block is $12d^{2}$, and $\val{p32.layers}$ of them are
  $\val{p32.blocks.pct}$ per cent of Program~\ref{prog:F01}'s
  $\val{p32.f01.params}$.}
\sumitem{33--35}{The cache is one key and one value per position per layer,
  which is $\val{p32.kv.per.token.mib}$ MiB a token and reproduces
  Program~\ref{prog:P03}'s $\val{p32.kv.gib}$ GiB. Double the sequence: the
  scores quadruple, the cache doubles.}
\sumitem{36--39}{Every derivation in this book about attention holds at
  initialisation, and five of them rest on hypotheses no trained model has
  been asked about. The list is printed rather than implied.}
\end{summarybox}

\begin{testexercises}
\item A block has width $d = 1024$. How many parameters are in its
  attention projections, and how many in its feed-forward block?
  \answerto{$4d^{2} = 4\,194\,304$ and $8d^{2} = 8\,388\,608$. The
  feed-forward block holds twice as many, at every width \dash{} the ratio
  does not depend on $d$.}
\item A model has $16$ layers of width $2048$. Estimate its parameter count
  from the blocks alone.
  \answerto{$12d^{2}L = 12 \times 2048^{2} \times 16$, which is about
  $805$ million. The embedding and the output head are on top of that.}
\item Two queries and two keys are produced by projections $W_Q$ and $W_K$.
  Which of $W_Q$, $W_K$ and $W_Q W_K\T$ can be recovered from the scores
  alone?
  \answerto{Only the product. The scores are $x\T W_Q W_K\T y$, so any pair
  of matrices with the same product gives identical scores \dash{} which is
  why the two are often stored and analysed as one matrix.}
\item A residual stack has $n$ layers and every $f'_k$ happens to be exactly
  $-1$. What is the gradient, and does the identity path save it?
  \answerto{$\prod(1 + (-1)) = 0$. The identity path contributes $1$ and the
  other $2^{n} - 1$ paths cancel it exactly. \enquote{One term is $1$} does
  not mean \enquote{the sum is at least $1$}, and this is the case that
  shows it.}
\item A serving budget allows $32$ GiB for the cache at the shapes of
  section~\ref{sec:P32-count}. How many $\val{p32.seq}$-token conversations
  fit?
  \answerto{Eight, since each is $\val{p32.kv.gib}$ GiB. Halve the context
  and sixteen fit \dash{} the cache is linear in the sequence, so this trade
  is one for one.}
\item Which quantity in a block grows quadratically with the sequence
  length, and which grows linearly?
  \answerto{The number of scores grows quadratically; the cache grows
  linearly. Both are properties of one layer, which is why a capacity plan
  has to name which one it has run out of.}
\end{testexercises}

\canyou

\begin{furtherproblems}
\item Show that $\Ex\lVert W_Q W_K\T\rVert_F^{2} = d_k$ when the entries of
  both matrices have variance $1/d$, by counting the terms.
  \answerto{$M$ has $d^{2}$ entries; each is a sum of $d_k$ products of two
  independent entries, so each has variance $d_k/d^{2}$; and
  $d^{2} \times d_k/d^{2} = d_k$. The stream's width cancels, which is why
  the divisor depends on the head width and not on the model's.}
\item A head is made narrower and the number of heads is raised to keep the
  total width fixed. What happens to the parameter count, and what happens
  to the rank of each head's $M$?
  \answerto{The parameter count is unchanged. Each $M$ has rank at most
  $d_k$, which falls, so the same parameters express a more constrained set
  of comparisons.}
\item Take the expansion $\prod_{k=1}^{n}(1 + f'_k)$ and write down the sum
  of the terms that pass through exactly one layer.
  \answerto{$\sum_k f'_k$. Together with the identity path's $1$ that is the
  first-order approximation $1 + \sum_k f'_k$, which is why a shallow
  residual stack behaves like a sum of its layers rather than a product.}
\item The scores are $n^{2}$ numbers and the cache is linear in $n$. At what
  sequence length does the score matrix, at two bytes a number and one head,
  match the cache of section~\ref{sec:P32-count}?
  \answerto{Set $2n^{2} = 2 \times \val{p32.layers} \times \val{p32.d}
  \times 2 \times n$, giving $n = 2 \times \val{p32.layers} \times
  \val{p32.d}$, which is $262\,144$ \dash{} far beyond the context lengths
  here, which is why the cache dominates at the sequences people serve and
  the scores dominate the arithmetic anyway.}
\item Program~\ref{prog:F12} shows a plain chain can also \emph{explode},
  and section~\ref{sec:P32-residual} says the residual path does not fix
  that. Why not?
  \answerto{The identity path bounds the gradient from below, not above. The
  other paths are unchanged, so a stack whose factors are large still
  multiplies them \dash{} which is what the normalisation of
  section~\ref{sec:P32-norm} is for, and why the two appear together in
  every block.}
\end{furtherproblems}
